% Part: first-order-logic
% Chapter: introduction
% Section: syntax

\documentclass[../../../include/open-logic-section]{subfiles}

\begin{document}

\olfileid{fol}{int}{syn}

\olsection{النحو}

أول ما يلزم ضبطُ سلاسل الرموز التي تُعد من ال!!{sentence}s في
منطق الرتبة الأولى. نؤخر التعريف الدقيق إلى موضعه، ونمهد له هنا
بالأمثلة. معجم هذا المنطق قسمان. فأما أولهما فرُموز نعيّن بها أشياء
أو نحكم بها عليها: التعيين بال!!{constant}s، والحكم بال!!{predicate}s.
فمثلًا، نأخذ $\Obj a$ !!a{constant} يعيّن شيئًا واحدًا، ثم نحكم
عليه بال!!{sentence}~$\Atom{\Obj P}{\Obj a}$. ومتى تصورنا معنى
لـ«$\Obj a$» و«$\Obj P$»، أمكن أن نفهم
$\Atom{\Obj P}{\Obj a}$ فهم جملة من اللغة الطبيعية؛ وهذا ما
يتمرن عليه الطالب في دراسته الأولى للمنطق الصوري. ومن هذه
ال!!{sentence}s البسيطة نؤلف جملًا أعقد بالقسم الثاني من المعجم،
وهو الرموز المنطقية: الروابط والمكممات. ومن أمثلتها
$(\Atom{\Obj P}{\Obj a} \land \Atom{\Obj Q}{\Obj b})$
و~$\lexists[\Obj x][\Atom{\Obj P}{\Obj x}]$.

ولا تستقيم التعريفات الدقيقة للدلالة ولقواعد أنساق ال!!{derivation}
اللازمة للبحث الفوق منطقي إلا بعد ضبط ما نعنيه بــ!!a{sentence}
في منطق الرتبة الأولى. أصل الفكرة واضح: نؤلف
!!{sentence}s بسيطة من ال!!{predicate}s وال!!{constant}s، كالصيغة
~$\Atom{\Obj P}{\Obj a}$، ثم نزيدها تركيبًا بالروابط والمكممات.
لكن أي قواعد تجيز تأليف ال!!{sentence}s المركبة؟ إن لم نصرح
بها، لم يتم حدّ «!!{sentence} في منطق الرتبة الأولى».
ويجب في ذلك مراعاة أمور: أن تدخل السلاسل الصحيحة وحدها في
ال!!{sentence}s؛ وأن يوضع التعريف على وجه يتيح البرهان الرياضي
على خواص \emph{جميع} ال!!{sentence}s؛ وأن تُضبط العمليات الأولية
على ال!!{sentence}s، ومنها «إبدال كل $\Obj x$ في~$!A$
بـ~$\Obj b$». ولوجود المكممات وال!!{variable}s لا يتضح لأول وهلة
كيف تُجرى هذه الأعمال. أَنعد، مثلًا،
$\lexists[\Obj x][\Atom{\Obj P}{\Obj a}]$ !!a{sentence}؟
وماذا نقول في
$\lexists[\Obj x][\lexists[\Obj x][\Atom{\Obj P}{\Obj x}]]$؟
وما الذي ينتج إذا «أبدلنا $\Obj x$ بـ~$\Obj b$ في
$(\Atom{\Obj P}{\Obj x} \land
\lexists[\Obj x][\Atom{\Obj P}{\Obj x}])$»؟

\end{document}
